% =====================================================================
%  ANONYMISED SUBMISSION DOCUMENT
%  Radiology: Artificial Intelligence -- Original Research
%
%  Derived from the full manuscript by cutting. See README.md in this
%  directory for what was cut and why.
%
%  DOUBLE-ANONYMIZED REVIEW. Nothing in this file may identify the
%  authors: no institution, no funding, no repository URL, no archive
%  DOI, no software name, no file paths, no "our previous study". The
%  full (non-anonymised) title page is a SEPARATE file and is not part of
%  this document.
%
%  Order of this single document, per the journal's instructions:
%    Abbreviated Title Page -> Abstract -> Main Body -> Acknowledgments
%    -> References -> Figure Legends -> Tables (one per page).
%  Figures themselves go in their own combined document (figures.tex).
%
%  Format rules applied here: double-spaced, ragged right (not
%  right-justified), 11 pt Times New Roman, pages NOT numbered.
%
%  Build: tectonic main.tex   (bibliography is handled by tectonic)
% =====================================================================
\documentclass[11pt,letterpaper]{article}

\usepackage[T1]{fontenc}
\usepackage[utf8]{inputenc}
\usepackage{mathptmx}          % Times New Roman text and math
\usepackage[margin=1in]{geometry}
\usepackage{setspace}
\usepackage{amsmath}
\usepackage{graphicx}
\usepackage{array}
\usepackage{booktabs}
\usepackage{siunitx}
\usepackage[font=small,labelfont=bf]{caption}
\usepackage[numbers,sort&compress]{natbib}
\usepackage{microtype}
\usepackage[hidelinks]{hyperref}

\sisetup{group-separator={,},group-minimum-digits=4}

\doublespacing
\raggedright
\pagestyle{empty}              % the journal asks that pages not be numbered
\setlength{\parindent}{0pt}
\setlength{\parskip}{0.6em}
\renewcommand{\arraystretch}{1.1}

\begin{document}

% =====================================================================
%  ABBREVIATED TITLE PAGE  (anonymized)
%  Four required items: Title, Article Type, Summary Statement, Key Points.
% =====================================================================

{\large\bfseries What a Slice-Level Benchmark Certifies without the Pixels:
An Audit of Seven Public Imaging Datasets\par}

\vspace{0.8em}

\textbf{Article Type:} Original Research

\vspace{0.4em}

\textbf{Summary Statement}

\textbf{A model that never sees a pixel reached about half of the published
margin over chance on volumetric imaging benchmarks read slice by slice, and on
the largest benchmark the same scores carried no patient-level information at
fixed stack depth.}

\vspace{0.4em}

\textbf{Key Points}

\begin{itemize}\setlength{\itemsep}{0.1em}
\item On \num{18938} patients, one pixel-blind score vector read 0.737 (95\%
confidence interval: 0.735, 0.740) per slice and, after mean aggregation, 0.453
(95\% confidence interval: 0.445, 0.461) per patient; the sub-chance patient
reading is a stack-depth confound of the mean operator; in a stratified
re-analysis, within strata of equal stack depth, the reading is 0.497 (95\%
confidence interval: 0.487, 0.508) against 0.449 unstratified in that draw.
\item Against peer-reviewed comparators, a model with no access to the images
reached a median 0.469 (interquartile range: 0.437, 0.490) of the published
margin over chance over nine benchmark-arms; those nine arms are four
benchmarks, six of them being the six label columns of one cohort, and taking
one arm per benchmark the median is 0.474 (range, $-0.002$ to 0.490).
\item On one of those four benchmarks the same model reached none of the
published margin ($-0.002$), so a benchmark can register as clean on this
measure.
\end{itemize}

\vspace{0.4em}

\textbf{Abbreviations}

AUC = area under the receiver operating characteristic curve;
CI = confidence interval;
DWI = diffusion-weighted imaging;
ICH = intracranial hemorrhage;
IQR = interquartile range.

\vspace{0.4em}

\textbf{Keywords:} CT; MRI; Computer Applications---Detection/Diagnosis;
Technology Assessment; Brain/Brain Stem.

\clearpage

% =====================================================================
%  ABSTRACT  (structured, four sections, <=250 words)
% =====================================================================
\section*{Abstract}

\textbf{Purpose:} To quantify how much published slice-level performance on
volumetric imaging benchmarks is reachable from the benchmark's
label file with no images, and what those scores show per patient.

\textbf{Materials and Methods:} Retrospective secondary analysis of
de-identified public label files from seven imaging benchmarks released
2016--2024, including NYU fastMRI; analyses ran July 2026. The largest,
the 2019 RSNA Intracranial Hemorrhage benchmark, held \num{752802} slices from
\num{18938} patients; age and sex are not distributed. Five pixel-blind models (position, volume size, metadata, and
combinations), with a within-series permutation null, were fit on
subject-disjoint splits and read at both units. The share of a published margin over
chance reachable without images
(``trivial fraction'') was computed. Areas under the receiver operating
characteristic curve (AUCs) carry 95\% percentile bootstrap intervals
resampling subjects; no hypothesis test was performed.

\textbf{Results:} On the hemorrhage benchmark one pixel-blind vector read
AUC 0.737 (95\% CI: 0.735, 0.740) per slice and, after mean aggregation, 0.453
(95\% CI: 0.445, 0.461) per patient; in a stratified re-analysis at equal stack
depth, 0.497 (95\% CI: 0.487, 0.508). All six labels behaved
alike (gaps, 0.205--0.307); two full-cohort routes agreed to 0.003, four to 0.008.
Against peer-reviewed comparators the median trivial fraction was 0.469 (IQR:
0.437, 0.490) over nine arms of four benchmarks, or 0.474 per benchmark.

\textbf{Conclusion:} A pixel-blind model reached about half of the
peer-reviewed published margin over chance at the slice level; per patient the
same scores retained none of it, the sub-chance reading being a stack-depth
confound.

\clearpage

% =====================================================================
%  MAIN BODY
% =====================================================================
\section*{Introduction}

A three-dimensional acquisition is often labeled slice by slice---this slice
contains the lesion, that one does not---and performance is then
reported by pooling slices into a single area under the receiver operating
characteristic curve (AUC). The clinical question is almost always about a
patient. The gap is arithmetic: a slice-level
ranking can be dominated by \emph{where} a slice sits in the stack, because
findings are not uniformly distributed along an organ, and a score vector
encoding only that geometry ranks slices well and patients not at all.

None of this is new. Shortcut learning is named
\citep{geirhos2020shortcut}: Badgeley et al predicted hip fracture at AUC 0.78
from the radiograph, 0.91 with hospital process features, and
0.52---chance---on a test set balanced across those variables
\citep{badgeley2019hip}, and the mechanism recurs across
institutions, in COVID-19 radiograph models, and across thirteen datasets
\citep{zech2018pneumonia,degrave2021covid,ongly2024shortcut}. On the evaluation
unit, slice-level cross-validation inflated test accuracy
by 29\%--55\% and reached about 96\% on \emph{randomly labeled} data
\citep{yagis2021leakage,tampu2022leakage}, and more than half of surveyed
Alzheimer disease classification papers may have leaked
\citep{wen2020cnn,kapoor2023leakage}. Stronger baselines have been urged at
field level \citep{varoquaux2022failures,roberts2021pitfalls}.

The central construction used here is prior art. A public 2020 Kaggle notebook
on the RSNA-STR Pulmonary Embolism Detection challenge, ``Baseline with no
image,'' bins the label rate against relative slice position on the
training set and applies it to the test set with no pixels
\citep{osciiart2020baseline}. Yan et al, defining the
DeepLesion lesion-type task, publish a ``Baseline: Location feature'' row at
59.7\% eight-class accuracy against their own method's 90.5\%
\citep{yan2018deeplesiongraphs}.

Two questions remain open: that literature measures the
cost of a \emph{wrong}, leaky split, not what survives a \emph{correct},
patient-disjoint one, and it does not quantify how much of a reported
slice-level number a pixel-free model accounts for. The purpose of this study
was to quantify how much published slice-level performance on public volumetric
imaging benchmarks is reachable from published non-pixel fields under a
subject-disjoint split, and whether the same score vector, read per patient,
retains it.

\section*{Materials and Methods}

\subsection*{Study Design}

This was a retrospective secondary analysis of publicly released, de-identified
label files from imaging benchmarks released 2016--2024, against the numbers
their comparator papers report. Analyses ran in July 2026. No patient was
recruited. Institutional review board approval was not required: no data were
obtained through intervention or interaction with any living individual and no
identifiable private information was used. Participant age and sex are not
distributed with these label files and cannot be reported.

\subsection*{Benchmark Selection}

A dataset entered the audit if the four fields a positional null needs---subject
identifier, slice index, label, and train/test assignment---could be obtained
without pixel data. Seven benchmarks across eight label files
qualified: RSNA 2019 Intracranial Hemorrhage (ICH) \citep{flanders2020rsnaich},
fastMRI Prostate, T2-weighted and diffusion-weighted imaging (DWI) arms
\citep{tibrewala2023fastmriprostate}, DeepLesion \citep{yan2018deeplesion},
fastMRI+ knee \citep{zhao2022fastmriplus}, Duke Breast Cancer MRI
\citep{saha2021dukebreastmri,saha2018radiogenomics}, PI-CAI
\citep{saha2024picai}, and LUNA16 \citep{setio2017luna16}. RSNA ICH is the
largest: \num{752802} slices, \num{21744} series, \num{18938} patients.

Every RSNA ICH number here rests on a slice ordering its official release does
not contain: that release is keyed by instance identifier, with no subject
identifier and no slice index. Both came from a public, MIT-licensed, pixel-free tabular mirror of the same
release, joined to the official labels on the instance identifier; the join
matches the comparator's description of that file to one row (\num{21744}
series; \num{752802} rows against \num{752803}) \citep{burduja2020ich}. That the
mirror's within-series counter is anatomic order was tested on all \num{21744}
series, not assumed: hemorrhage is spatially contiguous, so a
$z$-ordered index makes positive slices form few runs. Observed 1.384 runs
per series against 7.167 expected under random placement, and
7.170 with labels permuted within series; a failure returns that random
expectation, as the same test does on another challenge's release, which was
excluded on that measurement. The test fixes order but not orientation, which
does not matter: the positional model is invariant to reversing the
axis.

% -------------------------------------------------------------------------
% DO NOT DELETE OR ABRIDGE THE FOLLOWING PARAGRAPH. It is the acknowledgement
% mandated by the NYU fastMRI Data Sharing Agreement, paragraph 6(c), which
% closes: "further acknowledge that inclusion of some variation of the
% language shown above is mandatory." Paragraph 6(a) additionally requires
% the literal string "NYU fastMRI" in the abstract; it is there.
% This paragraph names no author of this manuscript and does not breach the
% anonymisation requirement.
% -------------------------------------------------------------------------
Data used in the preparation of this article were obtained from the NYU fastMRI
Initiative database (\url{https://fastmri.med.nyu.edu})
\citep{knoll2020fastmri,zbontar2018fastmri}. As such, NYU fastMRI investigators
provided data but did not participate in analysis or writing of this report. A
listing of NYU fastMRI investigators, subject to updates, can be found at
\url{https://fastmri.med.nyu.edu}. The primary goal of fastMRI is to test
whether machine learning can aid in the reconstruction of medical images.

\subsection*{The Zero-Image Baseline Family}

Five null models were fit on training rows and scored test
rows: a constant \emph{prevalence} predictor, the chance anchor; a
\emph{positional} model, a 20-bin estimate of the label rate
given relative slice position; a \emph{volume-size} model, the number of slices
in the volume; a \emph{metadata} model, a depth-limited classification tree
over the label file's remaining non-image columns, acquisition
and administrative in every benchmark but PI-CAI, whose file also
supplies age and prostate-specific antigen level; and a \emph{combined}
position-plus-metadata tree. Relative position is
$(\text{slice}-\text{minimum})/(\text{maximum}-\text{minimum})$ within a volume,
making volumes of different depth comparable. The positional baseline was also
run over a 5/10/20/50 bin sweep and a \emph{no-fit} variant,
$-\lvert\text{relative position}-0.5\rvert$, which uses no training data;
a result surviving both cannot be a binning artifact. Columns that are the
label under another name, or derived from the images, were
excluded; every column is recorded either way.

\subsection*{Evaluation, Intervals, and Nulls}

Every baseline produces one score vector per test set, read at \emph{both}
units: slice-level AUC, and patient-level AUC after averaging each subject's
slice scores, with a maximum-aggregated reading alongside. Splits were subject-disjoint: a benchmark's published train/test assignment
where one exists, otherwise a five-fold split scored out of fold and pooled.
Confidence intervals (CIs) are percentile bootstrap over \emph{subjects}---a
subject drawn twice contributes all its slices twice---with 2000 replicates for
the primary hemorrhage analysis and 200, 1000, or 2000 elsewhere, stated per row
in the table notes. The naive slice-resampled interval was also computed and
labeled as the incorrect one; coverage of both was measured on simulated data
with a closed-form true AUC (0.688; 200 datasets of 20 patients $\times$ 15
slices). Each baseline was reported against its own permutation null; for the
positional model that null shuffles labels \emph{within each series}, which
preserves prevalence, subject clustering, and stack depth.

\subsection*{The Comparison Statistic}

\begin{equation*}
\text{trivial fraction}=\frac{\text{best zero-image baseline}-\text{chance}}
{\text{published}-\text{chance}}
\end{equation*}

with chance $=0.5$ for AUC and the majority-class rate for multiclass accuracy,
computed for every (benchmark, comparator) pair. Rows are not
independent---one comparator table supplies several systems, and one benchmark
several label columns---so the primary summary takes each benchmark-arm's
\emph{strongest} published system, the conservative choice because a stronger
comparator enlarges the denominator, and is read against the count of
\emph{benchmarks} rather than of arms. Rows with a comparator that had not been
peer reviewed are summarized separately and labeled as such. The
fraction is undefined when the published number is at or below chance; values
above 1 and below 0 are reported unclipped. Its interval propagates uncertainty
in the baseline only, so it is too narrow as a statement about the ratio.

\subsection*{Statistical Analysis}

This is an estimation study: no significance test was pre-specified and no $P$
values are reported. Every AUC carries a 95\% percentile bootstrap CI resampling
subjects, with replicate count and seed recorded per run and degenerate
replicates counted rather than dropped; one reported range is instead a
split-to-split spread and is labeled as such. Analyses used Python 3.14 (Python
Software Foundation, Wilmington, Del) with NumPy 2.4 and pandas 3.0, and, for
the reimplementation only, scikit-learn 1.8.

\section*{Results}

\subsection*{One Pixel-Blind Score Vector, Read at Two Units}

Every value here is computed from a published label file; no published number
enters. On the RSNA ICH official training file---whose
own official metric is per slice---a 20-bin estimate of the
label rate given relative slice position, fit out of fold on a
subject-disjoint five-fold split, produces \emph{one} score vector. Read at the slice level it reaches AUC 0.737 (95\% CI:
0.735, 0.740). Read at the patient level, after mean aggregation, the same
vector reaches 0.453 (95\% CI: 0.445, 0.461)---an interval lying wholly below
chance. Nothing changes between the readings except the unit; the difference is
0.284, and all six labels behave alike, gaps 0.205--0.307 (Table 1,
Fig 1A). Maximum aggregation is vacuous here: every volume spans all 20 bins,
so every patient takes the same maximum, the score takes one value per fold, and
the reading ranks fold identity alone (0.493; 95\% CI: 0.485, 0.501).

Three controls ran on the same path. The within-series permutation null sits at
0.502 per slice, so the excess is 0.236; per patient it sits \emph{above}
chance, at 0.523 here and 0.561 under the audit tool's folds, so the observed
0.453 sits 0.069 below what the estimator reaches with nothing to find. The
constant predictor is a measured floor rather than 0.500 once folds are pooled,
at 0.492 per slice (Table 1). The naive slice-resampled interval is 1.51--1.99
times too narrow across these six rows and undercovers in simulation (Fig 1B).

The two full-cohort routes, one a second implementation sharing no code
with the audit tool, agree to 0.003 AUC at both units; all four span
0.008 (Table 2). Three further checks, given with Table 1, show that the
slice-level reading is not a binning artifact, that the null model is not
overfitting, and that metadata adds nothing here.

The sub-chance patient reading is a stack-depth confound of the mean operator,
not a reversed positional signal. A patient's mean is one fixed 20-value curve sampled at that patient's slice
positions, so stack depth accounts for three-quarters of its variance (cohort
range 0.133--0.150); and depth tracks the label here---series with hemorrhage
average 33.9 slices against 35.1 without, and depth alone reads 0.402 (95\% CI:
0.394, 0.410) per patient, the direction the mean inherits. Held at fixed depth
the anti-ranking goes: 0.497 (95\% CI: 0.487, 0.508) within exact depth and
0.491 (95\% CI: 0.482, 0.501) within 5-slice strata, both covering 0.500
(Table 1). The same mechanism drives the patient-level bin sweep, which unlike
the slice-level one is not bin-robust (0.437, 0.445, 0.456, 0.632 over 5, 10,
20, 50 bins). Twenty bins is pre-specified.

Of the six other benchmark-arms in Fig 2, four show the same divergence and two
do not (Table 3). DeepLesion does not collapse and should not: its labels are
anatomic regions, so position predicts them at both units. All eight body-part
arms rank above chance at both (slice 0.691--0.977, patient 0.642--0.954); the
one carried into the tables, pelvis, is the strongest and was chosen after
seeing them, so all eight are reported (Table 3). LUNA16
does not diverge between the units (0.534 slice, 0.581 patient). On Duke Breast
the patient-level value is undefined: all 922 of 922 patients are positive.

\subsection*{How Much of a Published Margin Is Reachable with No Pixels}

Seven benchmarks on eight label files produced 24 (benchmark, comparator) rows,
reported as a distribution rather than a verdict count (Table 4, Fig 3).

Against peer-reviewed comparators, taking each arm's
strongest published system, the median share of the published margin over
chance a pixel-blind model reaches is 0.469 (IQR: 0.437, 0.490; range,
$-0.002$ to 0.613; $n=9$ arms). Those nine arms are four benchmarks. Six are
the six label columns of RSNA ICH---one cohort, one slice ordering, one
positional model, against two columns of one comparator table
\citep{burduja2020ich}---at 0.395, 0.437, 0.469, 0.490, 0.510, and 0.613,
``any'' being the union of the other five. Correlated rather than independent,
those six contain the median and both quartiles.
One arm per benchmark gives a median of 0.474: RSNA ICH 0.490,
DeepLesion 0.480, PI-CAI 0.467, LUNA16 $-0.002$. Those four sit on four metrics
against four chance anchors, so the median describes this set rather than
estimating a field-wide quantity.

One benchmark of the four does not fire. On LUNA16, the same
positional estimator scored out of fold on a
scan-disjoint split gives sensitivity 0.0006 at one false positive
per scan---\emph{below} a random-score reference of 0.0027---against a published
combined-solution sensitivity above 0.95, for a trivial fraction of $-0.002$;
the challenge's own competition performance metric, a different quantity, is
0.0020 for the same scores
\citep{setio2017luna16}. A single negative shows the measure can
fail to fire; its false-positive rate is not estimable at
this $n$. On PI-CAI the
\emph{positional} baseline is exactly 0.500 at every bin setting, the correct
registration of ``inapplicable'' (one row per case, no
slice index) rather than a computed result; that arm is nonetheless not a null
here, because its strongest zero-image model
reaches 0.692 (95\% CI: 0.626, 0.755), a trivial fraction of 0.467
against a published case-level 0.91 \citep{saha2024picai}. That model
is a tree over the four non-image columns in that file, two of
which---age and prostate-specific antigen level---are standard clinical
predictors of clinically significant prostate cancer rather than acquisition
fields, so the row is a clinical-variable baseline and supports no inference
about acquisition confounding. Refitting on the two genuine acquisition
fields is not possible from the retained output; read singly they reach
0.547 and 0.552 per patient against 0.636 and 0.638 for the clinical pair
(Table 3).

Zero peer-reviewed comparisons are matched. Every row at or above a trivial
fraction of 1 comes from one benchmark, fastMRI Prostate, whose
comparator is a preprint still unpublished two years after
posting \citep{rempe2024kspace}; those six rows are
a worked example, not evidence for a general claim. On that group's own label
files, split, and evaluation unit, the zero-image
positional baseline reaches slice AUC 0.851 (95\% CI: 0.816, 0.887) against a
published 0.861---a trivial fraction of 0.973 (95\% CI: 0.876,
1.073)---while the same vector read per patient falls to 0.424
(95\% CI: 0.298, 0.547). That is the DWI arm, the one those authors
state they use. The general claim that
trivial baselines match published performance on medical imaging benchmarks is
not supported by these data, and not against the peer-reviewed
literature at all.

\section*{Discussion}

The unit at which a three-dimensional benchmark is read decides what its number
means. On \num{18938} patients a single pixel-blind score vector read 0.737
per slice and 0.453 per patient; only the first is what a paper would print. Against
peer-reviewed comparators the median share of the published margin
reachable with no pixels was 0.469 (IQR: 0.437, 0.490) over nine arms of four
benchmarks, 0.474 per benchmark.

A trivial fraction is a statement about an \emph{evaluation protocol}: a
pixel-blind model reaches that much of a reported margin over chance under it.
It is not a decomposition---the two models may exploit the same shortcut,
different ones, or overlapping ones---and not a claim about any model's
internals, which no label-file analysis could support. Reported continuously it
constrains both directions: about half the peer-reviewed published margin is
reachable with no pixels, and no peer-reviewed row reaches 1. Four benchmarks
cannot support a rate for the field.

Badgeley et al's control is \emph{negative}---balance the confounders and
watch the model collapse to 0.52 \citep{badgeley2019hip}. This one
is \emph{positive}: fit a model on the confounder alone and see what it
reaches, which an auditor can run without the images. Against the
leakage literature the distinction is the split, theirs wrong and this one
correct and patient-disjoint
\citep{yagis2021leakage,tampu2022leakage}; against Yan et al it is
the source of the position, theirs from a body-part regressor on the
image and this one from a published column
\citep{yan2018deeplesiongraphs}.

For a radiologist reading an imaging AI paper, the operative question is what a
reported AUC was computed over. On the hemorrhage benchmark the slice-pooled
number stood 0.205 to 0.307 above the per-patient number from the same
scores. The per-patient number is the one aligned with a per-patient
decision---but that is a statement about these benchmarks and the papers
reporting them, since no device summary or regulatory filing was examined here,
and nothing in these data speaks to the unit at which any cleared product was
validated. Read the patient-level value as a null, not a reversal: mean
aggregation carries stack depth, an acquisition property, into a per-patient
score, and depth is associated with the label here. A vector worth 0.737
per slice is worth nothing per patient, and a slice-pooled number hides that
completely.

For benchmark publishers, three fields already in the released label
file---a subject identifier, a slice index or $z$ position, and the official
train/test assignment---make every one of these checks runnable by anyone, at
no cost; RSNA ICH shows what it costs not to. What this protocol
cannot do is detect shortcuts living in the pixels---scanner texture,
burned-in annotation, body-part framing---which need the images.

\textbf{Limitations.} The construction is prior art
\citep{osciiart2020baseline,yan2018deeplesiongraphs}; measured here is how much
published performance it accounts for,
under a correct split, at two units, with subject-clustered statistics. No
peer-reviewed comparison is matched, and the only matched rows rest on a
preprint comparator. The headline patient-level value is not bin-robust, unlike
the slice-level one: it runs 0.437 to 0.632 over 5 to 50 bins (Results).
Twenty bins is pre-specified, in a protocol that will be identified on
acceptance. Two audited benchmarks are not pixel-free in the same
clean sense: fastMRI+ needs slice counts from image file headers, covering
199 of 1173 roster volumes, and LUNA16's false-positive-reduction track is
conditioned on a candidate list produced by image-based detectors, so ``zero
image'' there means ``given the published candidate list.'' The
PI-CAI comparison is across cohorts---published values on a hidden
1000-case testing cohort, the baseline on the public 1500-case development
set---so a strict reading makes those rows non-comparable; they are scored
because the caveat runs against the null. Age and sex are not in the public label
files, so no subgroup analysis was possible. Finally, it
covers only benchmarks whose label fields were obtainable without a
data-use agreement covering pixels.

In summary, one pixel-blind score vector ranked slices well
above chance and patients not at all, and reached about half
the peer-reviewed published margin over chance. Three fields
a benchmark already publishes make that checkable by anyone, without the
images.

% =====================================================================
%  ACKNOWLEDGMENTS  (anonymized)
% =====================================================================
\clearpage
\section*{Acknowledgments}

Acknowledgments are withheld from this version because they would identify the
authors, and will be supplied on acceptance.

\textbf{Use of large language models.} Large language model assistance was used
both in drafting and editing this manuscript and in writing the analysis code.
No value reported here was generated by a language model: every number in the
Results is a deterministic output of the released code, run on a published label
file, and the flagship estimate was reproduced by an independent
reimplementation written separately from the primary implementation, the two
agreeing to 0.003 AUC at both units on the full cohort (all four routes in
Table 2 span 0.008). Readers are directed to the released code rather
than asked to take the provenance on trust. Tool name, version, manufacturer,
and dates of access are given in the cover letter and will be named in this
section on acceptance, where naming them no longer risks identifying the
authors.

\textbf{Data and code availability.} The audit tool, the per-benchmark output
payloads, the label-table preparation scripts, the second implementation
of the intracranial hemorrhage unit-collapse computation, and the assembler
that builds the trivial-fraction table are released under a Creative Commons
Attribution licence in a public repository and archived with a persistent
identifier; the repository name, its address, and the commit identifier for the
revision used are withheld from this version for anonymized review and will be
stated in the Materials and Methods on acceptance, and are available to the
editors on request. No benchmark's pixel data are redistributed. The
worked-example label files from the NYU fastMRI Initiative are available by
application to that initiative under its data sharing agreement, which forbids
onward distribution of the dataset or of files derived from it; all other label
files used here are public.

% =====================================================================
\clearpage
\bibliographystyle{unsrtnat}
\bibliography{refs}

% =====================================================================
%  FIGURE LEGENDS  (figures themselves are in the separate figure document)
% =====================================================================
\clearpage
\section*{Figure Legends}

\textbf{Figure 1.} One pixel-blind score vector, read at two units, on the RSNA
2019 Intracranial Hemorrhage benchmark (\num{752802} slices, \num{21744}
series, \num{18938} patients). \textbf{(A)} Slice-level area under the receiver
operating characteristic curve (AUC) joined to patient-level AUC for each of
the six official labels, with 95\% subject-clustered bootstrap intervals
(2000 replicates) and
the within-series label permutation null at the patient unit. Nothing changes
between the two readings of a pair except the unit at which the ranking is
performed; every patient-level point estimate lies below chance, and four of
the six intervals lie wholly below it, the epidural and intraventricular
intervals crossing 0.5. The sub-chance patient-level values are a stack-depth
confound of the mean operator rather than a reversed positional signal; with
depth held fixed the ``any'' reading is 0.497 (95\% CI: 0.487, 0.508) (Table 1).
\textbf{(B)} Width of the correct subject-clustered
interval divided by the width of the naive slice-resampled interval on the same
six labels: the naive interval is 1.51 to 1.99 times too narrow. In simulation
at a nominal 95\%, the naive interval covered 46.5\% (93 of 200 datasets)
against 91.5\% (183 of 200) for the subject-clustered interval, with a mean
width 3.18 times narrower.

\textbf{Figure 2.} The same pixel-blind score vector read at two units, for
every audited benchmark-arm. Each point is one arm: its slice-level area under
the receiver operating characteristic curve (AUC) against the patient-level AUC
of the identical score vector after mean aggregation, with 95\%
subject-clustered bootstrap intervals on both axes (replicate counts in
Table 3). The dashed diagonal is the
line on which the two units agree, and the grey stem below each point is the
distance that arm falls from it. Of the seven arms for which both units are
defined, five have a patient-level interval that reaches chance or lies below
it; the two that do not are DeepLesion, whose labels are anatomic regions that
relative slice position predicts at either unit, and LUNA16, which moves little
at either. The DeepLesion point is the pelvis arm, the strongest of eight
body-part arms computed and chosen after all eight had been seen; the other
seven fall at 0.691 to 0.890 per slice and 0.642 to 0.838 per patient, all
above chance at both units, and all eight are listed in Table 3. Values are for
the 20-bin positional baseline except PI-CAI, whose positional baseline is
exactly 0.500 and whose strongest zero-image baseline is a tree over that label
file's four non-image columns, two of them clinical---age and prostate-specific
antigen level---rather than acquisition fields. The two arms defined at only one unit are drawn in
the margins rather than in the field, so that a value that does not exist cannot
be read as a low one: PI-CAI's label file carries no slice index, and Duke
Breast's patient-level AUC is undefined because all 922 of its 922 patients are
positive.

\textbf{Figure 3.} Trivial fraction---the share of a published margin over
chance reached by a pixel-blind model---for the strongest published comparator
per benchmark-arm. Peer-reviewed median 0.469, IQR 0.437 to 0.490 over nine
arms; those nine arms come from only four benchmarks, six of them being the six
label columns of RSNA ICH, and one arm per benchmark gives a median of 0.474
(Table 4). The
shaded band is that IQR, and error bars are 95\% subject-clustered bootstrap
intervals propagating uncertainty in the baseline only, at the replicate counts
in Table 3. Rows whose comparator
is a preprint that has not been peer reviewed are marked as such; they are the
two highest here, at 0.973 and 0.981, and no row in this figure reaches 1.
LUNA16 is plotted at $-0.002$, computed from sensitivity at one false positive
per scan (0.0006) and not from the challenge's competition performance metric,
which for the same score vector is 0.0020; the reference of 0.0027, used in
place of 0.5, is the random-score value of that competition metric, the mean
over its seven false-positive operating points. No interval is available for
this row. The open circle on the PI-CAI row marks its
positional baseline alone, which is exactly 0.500 area under the curve---the
correct registration of ``inapplicable,'' because that label file has no slice
index---while the filled marker is its strongest zero-image baseline, a tree
over that file's four non-image columns, two of which (age and prostate-specific
antigen level) are clinical predictors rather than acquisition fields
(Table 3).

% =====================================================================
%  TABLES  (embedded, one per page)
% =====================================================================
\clearpage
\begin{singlespace}
\small

\textbf{Table 1: RSNA 2019 Intracranial Hemorrhage---One Pixel-Blind Score
Vector Read at Two Units, All Six Official Labels}

\vspace{0.5em}

\begin{tabular}{lccccc}
\toprule
Label & Slice prev. & Patient prev. & Slice AUC & Patient AUC (mean) & Gap\\
\midrule
Any hemorrhage   & 0.143 & 0.404 & 0.737 (0.735, 0.740) & 0.453 (0.445, 0.461) & 0.284\\
Epidural         & 0.004 & 0.016 & 0.712 (0.700, 0.725) & 0.492 (0.461, 0.524) & 0.220\\
Intraparenchymal & 0.048 & 0.247 & 0.751 (0.747, 0.755) & 0.480 (0.471, 0.490) & 0.271\\
Intraventricular & 0.035 & 0.174 & 0.805 (0.802, 0.808) & 0.497 (0.487, 0.508) & 0.307\\
Subarachnoid     & 0.047 & 0.182 & 0.690 (0.686, 0.695) & 0.485 (0.475, 0.496) & 0.205\\
Subdural         & 0.063 & 0.178 & 0.720 (0.717, 0.723) & 0.476 (0.466, 0.487) & 0.243\\
\bottomrule
\end{tabular}

\vspace{0.8em}

\textbf{Any hemorrhage, patient level, with stack depth held fixed}

\vspace{0.4em}

\begin{tabular}{lcc}
\toprule
Patient-level reading & AUC & Pairs compared\\
\midrule
Mean aggregation, unstratified            & 0.449 (0.441, 0.457) & \num{86360472}\\
\quad within exact stack depth (36 strata) & 0.497 (0.487, 0.508) & \num{10065308}\\
\quad within 5-slice depth strata          & 0.491 (0.482, 0.501) & \num{25685534}\\
Stack depth alone                          & 0.402 (0.394, 0.410) & \num{86360472}\\
Maximum aggregation                        & 0.493 (0.485, 0.501) & \num{86360472}\\
\bottomrule
\end{tabular}

\vspace{0.4em}
Note.---\num{752802} slices, \num{21744} series, \num{18938} patients. AUC =
area under the receiver operating characteristic curve; prev. = prevalence.
Data in parentheses are 95\% percentile intervals from 2000 bootstrap
replicates resampling patients, in both blocks. Nothing changes between the slice and patient
columns except the unit at which the ranking is performed. The within-series label
permutation null is 0.502 at the slice level and 0.523 at the patient level.
The patient-level null lies \emph{above} chance because permuting within a
series preserves stack depth, so the depth confound survives it; it is also
fold-dependent, running 0.523--0.561 over seven fold draws against 0.502--0.505
at the slice level. Each observed value is paired with the null drawn under the
same fold assignment: under the audit tool's folds the pair is 0.456 against
0.561, a gap of 0.104 rather than 0.069.
The constant predictor scores 0.492 per slice and 0.501 per patient. It is
0.500 by construction within a single test set but not once folds are pooled,
because each fold emits its own training prevalence, which falls as that fold's
own prevalence rises; the five fold prevalences reproduce the 0.492 in closed
form with no other input. Folds are balanced on the patient-level label, which
is why the patient reading is 0.501 and the slice reading is not. Across the six
labels the constant predictor runs 0.480 (epidural) to 0.495 (intraparenchymal)
per slice and 0.498 to 0.504 per patient; each label's value is the measured
floor for that label's reading.
Three further checks ran on the same full-cohort output. The slice-level
reading is not a binning artifact: 0.716, 0.733, 0.738, and 0.745 over 5, 10,
20, and 50 bins, with the fit-free centrality score at 0.735. The null model is
not itself overfitting: apparent training AUC 0.738 against held-out 0.738. And
metadata adds nothing here: the combined position-plus-metadata tree reaches
0.718 against its own permutation null of 0.718.
Stack depth alone reads 0.402 as a raw score and 0.591 once its direction is
fitted (the volume-size baseline); these are the same information read with
opposite sign, and mean aggregation inherits the first.
The lower block re-reads the same score vector with stack depth held fixed,
under an independently drawn fold assignment whose unstratified patient-level
AUC is 0.449; over six fold assignments the slice-level AUC varies by less than
0.001 and the patient-level runs 0.445--0.457. Stratified values are pooled
within-stratum concordance with ties counted as half, over the pairs shown;
stack depth is a patient's slices per volume. Maximum aggregation is degenerate:
every volume holds at least 20 slices and spans all 20 bins, so all \num{18938}
patients take one of only five scores, one per fold, and the six labels' spread
of 0.486--0.505 is the spread of that artifact.

\end{singlespace}

\clearpage
\begin{singlespace}
\small

\textbf{Table 2: Independent Routes to the Intracranial Hemorrhage Unit
Collapse}

\footnotesize

\vspace{0.5em}

\setlength{\tabcolsep}{3.5pt}
\begin{tabular}{>{\raggedright\arraybackslash}p{0.29\textwidth}cccc}
\toprule
Route & Cohort & Reps & Slice AUC & Patient AUC\\
\midrule
Audit tool, 5-fold subject cross-validation & 1500-patient subsample & 200 & 0.731 (0.723, 0.739) & 0.462 (0.431, 0.491)\\
Independent implementation, different folds & Same subsample & 2000 & 0.731 (0.723, 0.739) & 0.458 (0.428, 0.486)\\
Independent implementation & All \num{18938} & 2000 & 0.737 (0.735, 0.740) & 0.453 (0.445, 0.461)\\
Audit tool & All \num{18938} & 1000 & 0.738 (0.735, 0.740) & 0.456 (0.448, 0.464)\\
Split-geometry replication & All \num{18938} & --- & 0.738 [0.727, 0.750] & ---\\
\bottomrule
\end{tabular}

\vspace{0.4em}
Note.---AUC = area under the receiver operating characteristic curve; Reps =
bootstrap replicates. Data in round parentheses are 95\% percentile bootstrap
intervals resampling patients, at the replicate count shown. The last line is
\emph{not} a bootstrap and its bracketed range is not a CI: it is the
2.5th--97.5th percentile of the estimate over 200 re-draws of the comparator's
own published held-out geometry (744 series), that is, split-to-split spread
rather than sampling error, and it yields no patient-level value.
The two full-cohort routes agree to 0.003 AUC at both units (below 0.001 per
slice, 0.003 per patient); across all four routes the maximum disagreement is
0.007 per slice and 0.008 per patient. Disagreements are computed before rounding.
Values are printed to three decimals throughout, the most the smallest
replicate count here supports. The subsample row of the independent
implementation was recomputed at 2000 replicates for this table; its point
estimates are unchanged and its earlier recorded interval, (0.420, 0.486), came
from 100 replicates. The
second implementation shares no code with the audit tool: its own fold
assignment and seed, its own binning, its own AUC routine, its own clustered
bootstrap.

\end{singlespace}

\clearpage
\begin{singlespace}
\small

\textbf{Table 3: Slice Level versus Patient Level for the Same Pixel-Blind
Score Vector, Every Audited Benchmark-Arm with Both Axes}

\vspace{0.5em}

\begin{tabular}{lcc}
\toprule
Benchmark-arm & Slice AUC & Patient AUC\\
\midrule
RSNA ICH, any hemorrhage & 0.737 (0.735, 0.740) & 0.453 (0.445, 0.461)\\
fastMRI Prostate, T2-weighted & 0.854 (0.812, 0.891) & 0.506 (0.381, 0.632)\\
fastMRI Prostate, DWI & 0.851 (0.816, 0.887) & 0.424 (0.298, 0.547)\\
fastMRI+ knee, meniscus tear & 0.873 (0.858, 0.886) & 0.510 (0.428, 0.592)\\
fastMRI+ knee, any annotated finding & 0.801 (0.779, 0.824) & 0.558 (0.470, 0.648)\\
Duke Breast, owner-defined slice task & 0.823 (0.811, 0.834) & Undefined (922 of 922 positive)\\
DeepLesion, pelvis vs rest & 0.977 (0.969, 0.984) & 0.954 (0.939, 0.967)\\
\quad mediastinum vs rest & 0.890 (0.871, 0.906) & 0.814 (0.782, 0.844)\\
\quad abdomen vs rest & 0.886 (0.868, 0.903) & 0.807 (0.774, 0.841)\\
\quad kidney vs rest & 0.877 (0.859, 0.894) & 0.838 (0.802, 0.870)\\
\quad liver vs rest & 0.876 (0.860, 0.890) & 0.777 (0.741, 0.812)\\
\quad lung vs rest & 0.872 (0.852, 0.892) & 0.803 (0.768, 0.836)\\
\quad soft tissue vs rest & 0.831 (0.777, 0.878) & 0.811 (0.769, 0.850)\\
\quad bone vs rest & 0.691 (0.618, 0.763) & 0.642 (0.555, 0.726)\\
LUNA16 candidates & 0.534 (0.513, 0.558) & 0.581 (0.538, 0.613)\\
PI-CAI, case level & Not applicable (no slice index) & 0.692 (0.626, 0.755), metadata only\\
\bottomrule
\end{tabular}

\vspace{0.4em}
Note.---AUC = area under the receiver operating characteristic curve; DWI =
diffusion-weighted imaging; ICH = intracranial hemorrhage. Data in parentheses
are 95\% percentile bootstrap intervals resampling subjects, from 2000
replicates except Duke Breast (1000) and LUNA16 (200). Values are for the
20-bin positional baseline except the PI-CAI row, whose positional baseline is
exactly 0.500. Two benchmarks, DeepLesion and LUNA16, do not collapse between
the units.
\emph{DeepLesion.} Eight body-part arms were computed and all eight are listed.
The pelvis arm is the one carried into the text, Figure 2, and the
trivial-fraction summary; it is the strongest of the eight at both units and was
selected after all eight had been computed, not on a rule stated in advance. The
other seven span 0.691--0.890 per slice and 0.642--0.838 per patient, and every
one of them lies above chance at both units, so the choice of arm affects the
illustration and not the conclusion drawn from it. All eight share one split,
one score construction, and the same 665 patients, and are therefore not
independent of one another.
\emph{PI-CAI.} That row's baseline is a depth-3 tree over the four non-image
columns of its label file---patient age, prostate-specific antigen level,
center, and scan year---and the fitted tree splits on all four. It is a
clinical-variable baseline, not an acquisition-metadata one, and the manuscript
draws no acquisition-confounding inference from it. Refitting it on center and
scan year alone is not possible from the retained output. Read one column at a
time on this split, patient-level AUC is 0.636 for age, 0.638 for
prostate-specific antigen level, 0.552 for scan year, and 0.547 for center;
scored instead out of fold over all \num{1476} development cases, the same four
give 0.600, 0.592, 0.517, and 0.504. The two clinical columns carry the row.
\emph{fastMRI Prostate.} On the DWI arm the fit-free centrality score reaches
slice AUC 0.841, so that row does not depend on fitting anything.
Rows scored on a benchmark's own published train/test split carry a constant
predictor of exactly 0.500: fastMRI Prostate (both arms), DeepLesion, and
PI-CAI. Rows scored out of fold and pooled do not, and their constant-predictor
value is the floor for that row: RSNA ICH 0.492, fastMRI+ knee meniscus tear
0.455, fastMRI+ knee any finding 0.483, Duke Breast 0.483, LUNA16 0.483. The
largest of these deviations, 0.045 on fastMRI+ knee meniscus tear, is of the
same order as that row's patient-level departure from chance, and the audit
tool raises this as a protocol warning on each of those runs.

\end{singlespace}

\clearpage
\begin{singlespace}
\small

\textbf{Table 4: The Trivial Fraction as a Distribution over 24 (Benchmark,
Published Comparator) Rows}

\vspace{0.5em}

\begin{tabular}{lcccccc}
\toprule
Set of rows & $n$ & Min & Q1 & Median & Q3 & Max\\
\midrule
Peer-reviewed, strongest per benchmark-arm & 9 & $-0.002$ & 0.437 & 0.469 & 0.490 & 0.613\\
Peer-reviewed, one arm per \emph{benchmark} & 4 & $-0.002$ & 0.350 & 0.474 & 0.482 & 0.490\\
Peer-reviewed, all rows & 18 & $-0.002$ & 0.455 & 0.485 & 0.514 & 0.889\\
Strongest per benchmark-arm, any comparator & 11 & $-0.002$ & 0.452 & 0.480 & 0.562 & 0.981\\
All rows & 24 & $-0.002$ & 0.469 & 0.512 & 0.910 & 1.655\\
Preprint comparator only (worked example) & 6 & 0.973 & 1.020 & 1.142 & 1.518 & 1.655\\
\bottomrule
\end{tabular}

\vspace{0.4em}
Note.---Trivial fraction $=$ (best zero-image baseline $-$ chance) $/$
(published $-$ chance). Q1 and Q3 are the first and third quartiles. The first
line is the primary summary: each benchmark-arm's strongest peer-reviewed
published system, which is the conservative choice because a stronger
comparator enlarges the denominator. Rows are not independent, since one
paper's table can supply several comparator systems for one benchmark, and one
benchmark several label columns.
\emph{The primary set of nine arms is not nine benchmarks; it is four.} RSNA ICH
supplies six of the nine, one per official label column, all scored on one
cohort with one slice ordering by one positional model against two columns of
one comparator table, and its ``any'' label is by construction the union of the
other five. Those six are strongly correlated rather than independent evidence,
they span 0.395 to 0.613, and they contain the median (subdural), the first
quartile (epidural), and the third quartile (any) of the nine. The second line
takes one arm per benchmark instead---RSNA ICH ``any'' 0.490, DeepLesion 0.480,
PI-CAI 0.467, LUNA16 $-0.002$---for a median of 0.474; substituting the median
of the six RSNA ICH labels (0.480) for ``any'' also gives 0.474. The quartiles
on that line are printed for consistency of format and should not be read as a
spread over four values. DeepLesion's row here is the 8-class lesion-type task,
a different task from the body-part arms carried under the same benchmark name
in Table 3. Those four benchmarks are also scored on four different metrics
against four different chance anchors---slice-level AUC, 8-class accuracy against
0.2361, case-level AUC, and sensitivity at one false positive per scan against
0.0027---so a median across them summarizes this set and is not an estimate on a
common scale. LUNA16's $-0.002$ is computed from sensitivity at one false positive per
scan (0.0006), \emph{not} from the challenge's competition performance metric,
which for the same score vector is 0.0020; the two are different quantities that
round alike. The reference of 0.0027, used in place of 0.5, is the random-score
value of the competition metric, that is, the mean over its seven
false-positive operating points; against the random-score value at the single
operating point actually used (0.0012) the fraction would be $-0.001$. Either
way this row is the only one in the primary set on this metric and this chance
anchor. Where a fraction carries an interval, that interval propagates
uncertainty in the baseline only, at the bootstrap replicate count listed for
that benchmark-arm in Table 3. No
peer-reviewed comparison is matched; the only rows
reaching 1 have a preprint comparator and are reported as a worked example.

\end{singlespace}

\end{document}
